\begingroup
\XeTeXlinebreaklocale "zh"
\XeTeXlinebreakskip=0pt plus .15em
\ChineseBody\fontsize{11}{19}\selectfont
\setlength{\parindent}{2em}\setlength{\parskip}{3pt plus .3pt}
\thispagestyle{opening}
\pdfbookmark[0]{作者自序：从开源社区到开源评价学}{origins.zh}
\phantomsection\label{origins:zh}
\markboth{Open Source Evaluatology}{Author\textquotesingle s Note · Chinese Original}
\noindent{\ChineseSans\fontsize{10}{14}\selectfont\color{FrameworkTeal}写在前面\quad /\quad 中文原文}\par
\vspace{10pt}
\noindent{\ChineseSans\fontsize{23}{31}\selectfont\color{HeadingInk}从开源社区到开源评价学}\par
\vspace{17pt}
% Paragraph 01
我开始使用开源软件，远早于我开始理解开源社区。本科时，Linux、Apache、Tomcat 就已经出现在我的学习和实践中。那时，我最直接的感受是这些软件好用、开放，遇到问题可以查文档、读代码。至于谁在维护它们，陌生人如何合作，一个项目为什么能够延续多年，我还没有认真想过。

% Paragraph 02
2017年的中国开源年会，让这些熟悉的软件名字有了另一层意义。我第一次真正接触到项目背后的人和社区。有人写代码，有人组织协作，有人帮助新人，也有人长期承担很少被外界注意的维护工作。开源开始从一组软件名称，变成我能够感受到的人与人之间的联系。后来回望，这次相遇是我走向开源研究的一个重要起点。

% Paragraph 03
此后，由于个人原因，我调整了工作，来到华东师范大学数据科学与工程学院。新的环境也带来了一个很实际的问题：在一个数据学院里，怎样开展开源教学和科研？兴趣需要找到可以长期投入的方向，社区中的观察也需要转化为能够研究、能够验证的问题。我逐渐意识到，开源协作留下的大量公开记录，恰好为数据科学提供了一个丰富而真实的研究场景。

% Paragraph 04
赵生宇的加入，是这段探索中很重要的一步。我们相识于开源社区。他从阿里巴巴开源办公室离开，得知我正在组建开源研究团队，便来到我这里攻读博士。我们有共同的兴趣，也带着不同的经验：社区实践让研究问题有了来处，学术训练则要求我们把问题说清楚，把证据做扎实。团队就这样一点点组织起来。

% Paragraph 05
2020年，我们发布了《GitHub 2019 数字年报》。围绕2019年约5.46亿条 GitHub 公开事件日志，我们尝试描绘开发者、项目和组织的活动，也观察中国开源参与者在全球协作中的位置。对我们而言，这是第一次比较系统地用数据呈现开源生态。报告获得的关注和肯定，给了这个年轻团队很大的鼓励，也让我们看到，社区和产业确实需要这样的观察能力。同年，生宇还发表了《开源点亮未来》一文。“开源点亮未来”也由此成为我们实验室一直沿用的口号。

% Paragraph 06
随着报告传播，我们也开始考虑，怎样把一次分析积累下来的能力保留下来。数据怎样采集和整理，指标怎样定义，计算过程怎样复现，结果怎样持续更新，这些工作都值得开放出来。我们将制作报告的过程与方法逐步整理成一个开源项目，这便是 OpenDigger 的由来。从一份报告走向可以共同维护的数据基础设施，我们对开源世界的研究也有了能够持续生长的依托。

% Paragraph 07
2021年，OpenDigger 进入中国电子技术标准化研究院推动的木兰开源社区孵化，随后逐步为国家有关部门和机构提供开源数据服务。它的使用场景也不断扩大，陆续支持了开源社的《中国开源年度报告》、思否的《中国开源码力榜》、中国开源软件推进联盟的《中国开源发展蓝皮书》，以及 InfoQ 的开源发展研究分析。每一次合作，都使我们更具体地理解：不同使用者需要的证据，并不完全相同。

% Paragraph 08
围绕 OpenDigger，我们与合作伙伴逐渐建设起一组生态项目：OpenLeaderboard 帮助人们浏览和比较，HyperCRX 把数据带到日常使用的代码托管界面，OpenGalaxy 让协作关系有了直观的空间形态；与开放原子开源基金会共建的 OpenSODA 大赛、与蚂蚁共建的 OSGraph、与湖北大学共建的 OpenTalent，则把这些能力带入数据创新、关系探索和人才培养等场景。许多新的研究问题，正是从这些具体使用中长出来的。

% Paragraph 09
工具越是被使用，我们就越不能回避指标背后的问题。一个人做了多少事，与这些事产生了怎样的影响，究竟是什么关系？某个项目的活跃度下降，应当怎样理解？一个排名能够帮助谁，又可能遗漏谁？这些问题很难单靠增加数据量来回答。活跃度研究、OpenRank 算法家族和后续实证工作，使我们不断向前推进，也不断提醒我们审视测量的含义和边界。

% Paragraph 10
也正是在这样的探索中，我与詹剑锋教授的交流逐渐深入。我们最早相识于他发起的 BenchCouncil 系列活动。2022年11月，我们举办了首届国际开源软件生态基准论坛，此后每年都深入参与 Bench 大会。开源研究中的数据、指标和应用问题，在这里与基准、测量和评价方法的讨论交汇起来。

% Paragraph 11
詹教授对评价学的探索，让我重新审视我们已经做过的工作。过去，我们往往从一项数据或一个应用出发，寻找合适的方法；评价学则促使我们进一步说明，评价的对象是什么，在什么条件下比较，关注谁所受到的影响，又依据什么作出价值判断。那些分散在算法、平台和案例中的问题，开始有了可以贯通起来的理论线索。

% Paragraph 12
2024年12月，在广州举行的 Bench 大会上，我作了关于开源评价学的主题报告，华东师范大学也获授首个国际（开源）评价科学与工程研究中心。对我们而言，这是多年实践与评价学理论进一步结合的一个节点。开源评价学由此有了更明确的研究定位：以评价学为理论基础，系统研究开放协作中的贡献、影响、价值及其评价方法。

% Paragraph 13
我尤其看重其中的双向关系。理论帮助我们发现原先没有说清楚的假设，实践则让我们知道，哪些概念需要细化，哪些方法需要修正。比如，社区中长期的维护和协调工作，是否在指标中得到恰当体现；同一套评价用于不同组织时，结果还能否作相同解释；评价公布后，人们的行为又会发生怎样的变化。把这些问题带回理论，才有可能让领域研究与基础研究相互促进。

% Paragraph 14
2025年 Bench 大会前后，《Evaluatology: The Science of Uncovering the Effects》英文专著正式发布，中文版也于近期面世。在詹教授的持续推动下，我看到 BenchCouncil 从会议交流逐步发展出期刊、专著、基准平台等相互支撑的学术生态。一个方向要走得长远，需要有人提出问题，也需要可以积累知识、讨论分歧和检验方法的共同空间。这给了我们很大的启发。

% Paragraph 15
这本《开源评价学：理论、方法与实证》，就是在这样的背景下逐步成形的。写作使我们重新整理多年的论文、博士学位论文、工具和应用案例，也要求我们回答一个更难的问题：它们为什么属于同一个研究方向？从理论基础到测量方法，从实证研究到智库应用与学科建设，全书四个部分尝试给出连贯的回答。我们希望读者既能理解一项方法怎样工作，也能判断它适合回答什么问题。

% Paragraph 16
从实践中总结理论，是一项需要克制的工作。我们自己的案例可以提供证据，却不能替所有开源社区发言；某项算法在已有数据上的有效性，也不能自动延伸为对一切情境的解释。写这本书，是为了把依据、方法和局限交代得更清楚，让后来者能够继续检验和推进。开源评价学的体系，需要在这样的积累中逐步站稳。

% Paragraph 17
共同完成本书的，除了生宇，还有我的两位博士生韩凡宇和夏小雅。小雅已经博士毕业，进入蚂蚁开源办公室工作，继续活跃在国际开源社区中；凡宇也预计于2027年完成博士学业。从社区相识，到共同研究，再到把长期积累写成一本书，这几位年轻人的成长始终与团队的发展交织在一起。能够与他们共同署名，是我作为老师很珍惜的一件事。

% Paragraph 18
当然，书上四个作者的名字，远不能代表全部参与者。实验室同学对数据、代码、实验和文稿的投入，学院同事的支持，业界伙伴与社区同仁提供的问题、反馈和合作机会，都是这项工作得以延续的条件。很多贡献并没有署名在封面上，却实实在在地留在了书中。谨在这里向所有同行者致谢。

% Paragraph 19
回望这些年，开源也帮助我回答了刚到数据学院时的那个问题。真实场景提出需求，应用推动方法创新，创新再回到场景中接受检验，并支持更长远的发展。“场景驱动应用、应用驱动创新、创新驱动发展”，在开源研究中有了具体的内容：一份报告可以催生一个数据项目，一个工具可以进入课程和社区，一个实践中的疑问又可以成为新的研究课题。教学、科研与服务因此有了相互连接的路径。

% Paragraph 20
我们由此为开发者、企业、社区、基金会和政府机构提供了数据与评价服务，也在这个过程中培养了约三十位青年人才。他们走向不同的岗位，有人继续做研究，有人进入产业与社区。对一个大学实验室来说，能够留下可被使用的成果，也能够培养愿意参与公共协作的人，是很实在的收获。

% Paragraph 21
这些探索也发生在中国开源不断发展的背景下。开源已经写入“十四五”和“十五五”规划纲要，从支持开源社区到推进开源体系建设，其公共意义日益清晰。以 DeepSeek（深度求索）等为代表的开放实践，又让更多人看到，中国的技术团队可以通过分享成果，吸引全球开发者参与，建立跨越地域的合作联系。作为研究者，我们有机会参与其中，也有责任帮助人们更准确地理解这些变化。

% Paragraph 22
2026年7月9日，在瑞士日内瓦举行的“人工智能向善”（AI for Good）全球峰会边会上，我代表中国科协联合国咨商开源创新专委会（CCOSI），参与发布《2026全球开源发展报告》（预览版）。报告引起的关注与交流，让我感受到，开源数据研究已经可以服务于更广泛的国际对话。从2019年 GitHub 数据的整理，到面向全球开源发展的报告，研究的范围在扩大，最初想要理解协作的愿望一直没有变。

% Paragraph 23
全球对话也让评价的责任更加具体。不同地区的参与条件不同，不同社区留下的数据也不一样。一个看似整齐的比较，可能掩盖重要的差异。我们希望科学的评价能够帮助人们理解贡献、识别需要支持的地方，并为合作提供可信的依据。让开放成果更容易被使用，让长期维护和公共贡献得到更恰当的认识，是这项研究值得继续做下去的理由。

% Paragraph 24
今天，人工智能的发展使这些问题更加迫切。更多人能够获得技术，并不意味着每个人都拥有同样的参与机会。如何让不同背景的开发者进入协作，如何让有价值的公共成果得到持续支持，如何让技术进步惠及更广泛的人群，仍然需要长期工作。我期待开源评价学能在其中贡献一种扎实的研究方式，使关于开放与普惠的愿望，有可以讨论、可以检验、可以改进的依据。

% Paragraph 25
这本书面向国际读者，以英文写作。我仍想用中文留下这段缘起，并附上英文译文。中文是我们几位作者共同的母语，也承载着我们理解世界、交流想法和表达感谢的许多细节。保留它，是希望读者在概念、模型和图表之外，也能看见这项工作从什么地方出发，与哪些人相遇，又为什么坚持到了今天。

% Paragraph 26
从最初使用开源软件，到走进社区，再到与同行共同探索开源评价学，一路走来，最重要的始终是人与人的相遇，以及由此生长出的共同工作。愿这本书成为其中的一份贡献。也愿更多读者从这里接着往前走，让开放协作的价值被更好地理解，让愿意参与的人得到更多机会，让我们共同依赖的开源生态继续生长。

\par\vspace{12pt}\begin{flushright}\ChineseBody 王伟\\2026年11月，上海\end{flushright}
\clearpage\endgroup
